\section{The tuning confound}
\label{sec:tuning}

Everything in Section~\ref{sec:results} is measured at one learning rate: the
pinned recipe's $0.0226$, applied unchanged to every arm. That is the right way
to isolate a geometry change, and it is the wrong way to compare optimizers. An
optimizer modification that merely lets the step size grow will look like a
geometry win at any fixed rate, and the only way to tell the two apart is to
tune each arm separately. This section reports what happened when we tried, and
the answer reframes the rest of the paper.

\subsection{On a short screening schedule, the optimum shifts three times higher}

We screened seven matrix learning rates --- the recipe value scaled by
$2^{-1/2}$ through $2^{5/2}$ --- for four families: the baseline and three
diagonal output-side variants. Each of the 28 trials is a self-contained
320-update, 167{,}772{,}160-token run with its own complete 20\%-flat,
80\%-linear-decay schedule, not an early stop of the full budget. Every family
received an identical token allowance, and ranking used the reserved held-out
selection loss only, with C4 held back as a reporting metric.

Figure~\ref{fig:lr} (left) shows the result. On this 6\%-budget schedule, the
baseline's loss falls from $4.566$ at the lowest rate to $\lrMuontBest{}$ at
$\lrMuontBestLR{}$ before rising again --- a short-horizon swing of
\lrHeadroom{} nats. Our first screening round, covering only the three rates
closest to the recipe, selected the highest rate it tested for all four
families; so did the second round, two rates higher. Only the seven-point round
produced interior optima on the short schedule.

\begin{finding}[Learning-rate sensitivity dominates short-schedule differences,
requiring full-budget tuning]
On the 320-update schedule, moving the baseline's matrix learning rate shifts
held-out selection loss by \lrHeadroom{} nats --- roughly \lrHeadroomRatio{}
times the \mlpdiagselmag{}-nat fixed-rate effect at full budget. Although
Subsection~8.3 shows that $\lrMuontBestLR{}$ over-steps the 5{,}120-update
budget by \waveTwelveMuontSelDiffMag{} nats, the scale of learning-rate
curvature makes independent full-budget tuning around $0.0226$ mandatory.
\end{finding}

\begin{figure}[t]
\centering
\begin{tikzpicture}
\begin{groupplot}[
  group style={group size=2 by 1, horizontal sep=1.7cm},
  width=0.46\textwidth, height=6.2cm,
  grid=major, grid style={gray!18},
  axis line style={gray!55}, tick align=outside,
  label style={font=\small}, tick label style={font=\footnotesize},
  every axis plot/.append style={line width=1pt, mark size=1.9pt},
]
\nextgroupplot[
  xlabel={Matrix learning rate},
  ylabel={Final held-out selection loss},
  xmode=log, log basis x=2,
  xmin=0.014, xmax=0.145, ymin=4.20, ymax=4.62,
  xtick={0.016,0.02262741696,0.032,0.04525483392,0.064,0.09050966784,0.128},
  xticklabels={.016,.023,.032,.045,.064,.091,.128},
  xticklabel style={font=\tiny, rotate=40, anchor=east},
  legend style={at={(0.97,0.95)}, anchor=north east, font=\scriptsize,
                draw=gray!40, fill=white, fill opacity=0.92, text opacity=1},
  legend cell align=left,
  title={320-update screening}, title style={font=\small},
]
\addplot[color=black, mark=triangle*] table[x=lr, y=muont] {figs/lr_sweep.dat};
\addlegendentry{MuonT (baseline)}
\addplot[color=laPurpleLine, mark=*] table[x=lr, y=inner_outer_leftdiag] {figs/lr_sweep.dat};
\addlegendentry{all diagonal}
\addplot[color=laBlueLine, mark=square*] table[x=lr, y=outer_leftdiag] {figs/lr_sweep.dat};
\addlegendentry{outer diagonal}
\addplot[color=laOrangeLine, mark=diamond*] table[x=lr, y=outer_leftdiag_fused] {figs/lr_sweep.dat};
\addlegendentry{fused outer}
\draw[gray!55, dotted, line width=0.8pt]
  (axis cs:0.02262741696,4.21) -- (axis cs:0.02262741696,4.60);
\node[font=\tiny, text=gray!60!black, rotate=90, anchor=south east]
  at (axis cs:0.02262741696,4.595) {recipe rate};
\nextgroupplot[
  ybar, bar width=16pt,
  symbolic x coords={all diagonal,outer diagonal,fused outer},
  xtick=data, xticklabel style={font=\scriptsize, rotate=20, anchor=east},
  ylabel={$\Delta$ selection vs baseline},
  ymin=-0.016, ymax=0.068,
  scaled y ticks=false,
  yticklabel style={/pgf/number format/fixed, /pgf/number format/precision=2},
  nodes near coords={\pgfmathprintnumber[fixed, precision=3, print sign]{\pgfplotspointmeta}},
  every node near coord/.append style={font=\tiny},
  nodes near coords align={vertical},
  title={5120-update transfer}, title style={font=\small},
]
% Generated by scripts/analyze.py. Do not edit.
\addplot[fill=crossred!65, draw=crossred!85] coordinates {(all diagonal,0.049294) (outer diagonal,0.050256) (fused outer,-0.002909)};

\draw[gray!60] (axis cs:all diagonal,0) -- (axis cs:fused outer,0);
\end{groupplot}
\end{tikzpicture}
\caption{The tuning confound in two panels. \textbf{Left:} held-out selection
loss against matrix learning rate on the 320-update screening schedule. All
four families improve steeply well past the recipe rate (dotted line) before
turning over; the baseline alone gains \lrHeadroom{} nats. Our first two
screening rounds both selected the top of the grid they tested, which is how a
tuning study tells you it has not found an optimum. \textbf{Right:} what
happened when the interior screening optima were carried to the full
5{,}120-update budget. Two of three geometries went from beating the baseline
to losing by about $0.05$ nats --- an order of magnitude larger, and opposite in
sign, to every effect in Section~\ref{sec:results}. Short-schedule learning-rate
selection does not transfer.}
\label{fig:lr}
\end{figure}

\subsection{At their own best rates, the geometries still led --- on that schedule}

At the interior optima the output-side arms retained an advantage over the
independently tuned baseline: $-0.019633$ for all-diagonal and $-0.014299$ for
outer-diagonal, both at $\lrBestLROuterLeftdiag{}$, and $-0.009946$ for fused
outer at $\lrBestLROuterLeftdiagFused{}$. Taken alone, this is the paper's most
favourable evidence --- larger gains than anything in Table~\ref{tab:main}, and
obtained after giving the baseline its own best rate.

It does not survive contact with the full schedule, and it also disagrees with
the full schedule at the recipe rate. On the 320-update schedule at the recipe
rate, all three candidates are \emph{worse} than the baseline, by
\shortBaseInnerOuterLeftdiag{}, \shortBaseOuterLeftdiag{} and
\shortBaseOuterLeftdiagFused{} nats, exactly inverting their full-budget
ordering. Whatever these geometries do, it interacts with schedule length,
which means a short screening schedule is not a valid instrument for ranking
them.

\subsection{Why short-schedule transfer inverted: horizon shift and LR mismatch}

We then ran the three screening-selected rates at the full 5{,}120-update budget
against a same-wave baseline at its own screening-selected rate ($0.064$).
Figure~\ref{fig:lr} (right) shows the relative deltas, and inspecting the
\emph{absolute} losses exposes two structural reasons why the transfer failed:

\begin{enumerate}
\item \textbf{The 320-step rate ($0.064$) severely over-steps the 5{,}120-step
  horizon, even for the baseline.} Carrying \muont{} from the recipe rate
  $0.0226$ to its 320-step optimum $0.064$ \emph{degraded \muont{}'s own
  full-budget late selection loss by \waveTwelveMuontSelDiff{} nats}
  (\cselWaveOneTwoMuont{} in Wave~12 versus \cselWaveOneZeroMuont{} at $0.0226$
  in Wave~10) and worsened final C4 by \waveTwelveMuontCfourDiff{} nats. A
  6\%-budget run naturally tolerates a $3\times$ larger step size; at 5{,}120
  updates, $0.064$ is already well past the full-budget optimum, and the
  recipe rate $0.0226$ was near \muont{}'s true full-budget optimum all along.
\item \textbf{The $+0.050$-nat candidate regressions reflect an LR-grid split across
  equivalent equations.} In Appendix~\ref{app:arms}, \outerdiag{} and \fused{}
  implement the \emph{exact same mathematical update}, differing only by FP32
  product fusion ($1.4\times 10^{-7}$ relative error). On the flat top of the
  320-step curve (Figure~\ref{fig:lr}, left), evaluation noise caused \alldiag{}
  and \outerdiag{} to select $0.09051$ ($4\times$ recipe) while \fused{} and
  \muont{} selected $0.064$ ($2.83\times$ recipe). Because $0.064$ was already
  past the full-budget stability cliff, pushing two arms another $\sqrt{2}$
  higher to $0.09051$ produced $+0.049294$ and $+0.050256$-nat regressions
  against $0.064$, whereas \fused{} --- tested at the \emph{same} $0.064$ rate
  as Wave-12 \muont{} --- still beat it by $-0.002909$ selection and
  $-0.001806$ C4 nats. Moreover, our primary width-thresholded geometry
  \mlpdiag{} was never included in this short-schedule sweep.
\end{enumerate}

\begin{finding}[Short-schedule optima over-step the full budget; full-budget
tuning must bracket $0.0226$]
Carrying 320-update learning rates to 5{,}120 updates degraded \muont{} itself
by \waveTwelveMuontSelDiffMag{} nats relative to $0.0226$, and split two
mathematically identical outer-diagonal arms across $0.09051$ ($+0.050$ deficit)
and $0.064$ ($-0.0029$ gain). Short-horizon screening is therefore an invalid
proxy for full-budget tuning, and the full-budget optimum lies near $0.0226$,
not $0.064$.
\end{finding}

\subsection{What this does and does not establish}

Together with Section~\ref{sec:mech-graft}'s proof that \mlpdiag{} and
\alldiag{} preserve direction norms ($\mechMlpdiagGraftNorm{}$ and
$\mechLeftdiagGraftNorm{}$) rather than inflating effective step sizes, the
horizon analysis narrows what remains open considerably. The recipe rate
$0.0226$ is not $26\times$ detuned at full budget --- $0.064$ is worse by
\waveTwelveMuontSelDiffMag{} nats --- and when tested at matched rates ($0.0226$
in Table~\ref{tab:main}, or $0.064$ for \fused{} in Wave~12), output-side arms
consistently lead \muont{}. What remains unmeasured is the exact full-budget
tuning curve in the neighborhood \emph{around $0.0226$} (specifically
$\{0.016, 0.0226, 0.032, 0.045\}$) for \muont{}, \mlpdiag{}, and \alldiag{},
which Section~\ref{sec:program} places at the top of the final confirmation
list.
